\section{Discussion}

\subsection{DiagBench as a Style-Boundary Compatibility Test}

DiagBench should be read as a style-boundary compatibility test, not as a global engineering leaderboard. Each probe changes the trusted state, model action, and failure consequence: P1 asks for entry discipline, P2 for verifier-feedback closure, P3 for contaminated-state reset, and P4 for policy-authorized selection. Frontier models arrive at these boundaries with different response-control priors. An over-act prior can help search persistence but hurt triage. A continuity-biased state prior can help ordinary feedback closure but harm corrupted-state recovery. An evaluator-role prior can help policy ranking while limiting generation-oriented search. The non-monotonic rankings are therefore not a nuisance; they are the measurement target.

This framing provides a compact account of the main apparent puzzles. Gemini leads P2 because its bounded local-edit style fits verifier-feedback repair, but the same generation-oriented style is less suited to P4 policy ranking. model_C leads P4 because it behaves like a disciplined evaluator, but that same profile does not produce strong P2 search closure. model_E gains most from the P3 state-summary intervention because its raw-history behavior shows strong continuity bias. Reasoning-style models shift P3 more than P1, P2, or P4 because state skepticism and explicit replanning are specifically useful at the contaminated-state boundary, while the other stages reward different priors.

\subsection{What the Evidence Establishes}

The strongest empirical claims are direct. First, the 12 complete model runs dissociate across P1--P4 under a shared verifier, and no model leads more than one headline stage. Second, replacing raw corrupted history with verifier-authored state summaries improves P3 recovery and reduces cascade on a frozen 36-task audit. Third, log-derived profile indicators are strongly stage-specific: action prior tracks P1, edit style tracks P2, and preference execution tracks P4. Fourth, the circuit audit reproduces the same diagnostic pattern in a second closed-form engineering domain.

Some claims remain interpretive rather than isolated causal conclusions. The five response-control dimensions are useful diagnostic indicators, but we do not claim they are complete or factorially separable. The P3 intervention changes the prompt-facing state representation, so its gain is consistent with verifier-authored state isolation but may also include compression, denoising, or formatting effects. The controlled-prompt audit (Section~\ref{sec:experiments}) tests whether Tier~1 behavioral fingerprints are causally upstream of performance; the null result supports the interpretation that these fingerprints reflect durable prior-boundary compatibility rather than prompt artifacts, but it does not isolate the exact cognitive mechanism. The circuit audit is a construct-validity check, not a full second benchmark. These boundaries matter: the paper's contribution is a verifier-grounded diagnostic framework and evidence that it exposes stable workflow dissociation, not a universal behavioral ontology.

\subsection{Implications for Engineering-Agent Systems}

The deployment implication is to route by boundary rather than by aggregate model rank. Search should be gated by a triage controller with balanced action discipline. Repair should use a model or prompt configuration whose edit style performs bounded, feasibility-preserving updates under oracle feedback. Recovery should not consume raw corrupted history as neutral context; continuity-biased models need verifier-authored state summaries or reset mechanisms. Final selection should be separated from generation because the best feasible-set seeker is not necessarily the best policy-conditioned evaluator. In short, engineering-agent systems should be verifier-gated, stage-routed, and policy-separated.

This also clarifies the role of classical optimization baselines. CMA-ES on the VEH P2 bank is useful calibration, but it does not replace the benchmark because DiagBench measures more than objective search. P1 tests actionability, P3 tests recovery after state contamination, and P4 tests policy-conditioned selection. Even inside P2, the key question is not only whether an optimizer can eventually improve utility, but whether an agent can close under a small oracle budget while preserving feasibility.

\subsection{Limits and Next Steps}

The main suite is anchored in VEH, and the hardened circuit audit is intentionally small. The analytical oracle makes the benchmark fast and auditable, but it narrows physical scope relative to high-fidelity simulation. P1 is a matched certification stress stage rather than a source-disjoint literature holdout. Some top-line gaps should be read with intervals rather than rhetoric, especially P3 and P4 near ties. The SG-gap audit currently covers four frontier models, and the P3 intervention covers a stratified subset rather than the full bank.

The next scaling step is not simply more VEH tasks. The portable object is the verifier-grounded decision scaffold: triage before search, bounded repair under feedback, recovery from corrupted state, and policy-conditioned selection among feasible candidates. Scaling this scaffold to thermal, structural, fluid, and circuit domains would test generality while preserving the same control boundaries. A second direction is factorial probing: vary feedback format, history length, policy explicitness, and verifier-authored summaries to isolate action prior, edit style, feedback obedience, state trust, and preference execution more cleanly than the current benchmark can.
